\section{Results}
\label{sec:results}

Across datasets, ensemble MC dropout has the best average rank for the TDC metric, calibrated Brier
score, and calibrated interval score (Figure~\ref{fig:global-rank-holm}). The Friedman tests find
clear differences among the methods for all three measures. The follow-up comparisons qualify the
ranking: the group not separated from ensemble MC dropout contains all five ensemble pipelines for TDC performance.
For Brier score, only the deterministic MLP and Laplace are separated from the leader. For interval
score, the group not separated from the leader includes LightGBM, MC dropout, and all five ensemble
pipelines.

\begin{figure*}[t]
  \centering
  \includegraphics[width=\textwidth]{global_rank_rail.pdf}
  \caption{Mean ranks across datasets, with rank 1 best and the empirical leader outlined. Marker
  shape distinguishes LightGBM, single MLP, and MLP ensemble pipelines. ``Not significantly
  worse'' marks methods that the Holm-adjusted comparison does not distinguish from the leader;
  ``Significantly worse'' marks methods that it does. Adjustment is performed separately for each
  evaluation measure across all 55 method pairs.}
  \label{fig:global-rank-holm}
\end{figure*}

Appendix~Table~\ref{tab:global-rank-holm} gives every mean rank and adjusted $p$-value.

LightGBM has the most empirical wins for both the TDC metric and Brier score, yet its performance is
strongly endpoint dependent. Its TDC ranks range from first to last, producing a substantially worse
average rank than ensemble MC dropout. For interval score, LightGBM and ensemble post-hoc VBLL each
win three datasets, while ensemble MC dropout ranks between first and fourth on every regression
dataset.

\paragraph{What individual datasets can distinguish.}
Average ranks summarize consistency across the suite; comparisons on the original score scale show
which methods the repeated evaluations can distinguish on each endpoint. Figure~\ref{fig:per-dataset-status-tally}
makes the contrast clear. LightGBM is often the sole method not ruled out as best, but it is also
ruled out on many other endpoints. Ensemble MC dropout wins fewer datasets but is not ruled out as
best across a broader range of them. The evaluations find some method difference for all 44
dataset-score combinations.

\begin{figure*}[t]
  \centering
  \includegraphics[width=\textwidth]{per_dataset_status_stacked.pdf}
  \caption{Per-dataset classifications from Hsu's MCB procedure after averaging the five model
  seeds and five outer folds within each of ten chemical partitions, giving ten paired
  complete-partition blocks. ``Statistically best'' means that the empirical winner is the only
  method not ruled out by MCB. ``Empirical best'' marks other observed winners, and ``Not ruled
  out as best'' marks additional methods not excluded by the test. The initial test finds a method
  difference for all 44 dataset-score combinations; each is assessed separately
  at the 5\% level, with the coverage interpretation given in
  Appendix~\ref{app:dataset-scores}.}
  \label{fig:per-dataset-status-tally}
\end{figure*}
\FloatBarrier

Appendix~\ref{app:dataset-scores} shows the corresponding scores and decisions for every dataset.

\paragraph{Evaluation-design stability.}
With one model seed, a single held-out fold selects the same leader as the completed grid in only
37.2\% of TDC-metric, 20.4\% of Brier, and 39.2\% of interval-score subsets
(Figure~\ref{fig:rank-trajectories}). The apparent leader from one fold may therefore differ from the
leader obtained after averaging the full evaluation. At the equal cost of five scores, one complete
cross-validation reproduces the full-grid rank ordering more closely than averaging five model seeds
on one held-out fold. Repeated complete cross-validations narrow the rank variation, with Brier score
settling most slowly.

\begin{figure*}[p]
  \centering
  \includegraphics[width=\textwidth,height=0.82\textheight,keepaspectratio]
    {evaluation_design_rank_trajectories.pdf}
  \caption{Suite-wide mean ranks under smaller evaluation designs. Open circles and bars show the
  median, interquartile range, and 5th--95th percentile range over paired finite-grid subsets;
  diamonds show the completed grid. Solid branches add folds and partitions, while dashed branches
  average model seeds. Percentages give the subsets that recover the sole full-grid leader. These
  empirical ranges are not confidence intervals.}
  \label{fig:rank-trajectories}
\end{figure*}

\paragraph{Tuning with a single validation split.}
To test whether a single train--validation split is sufficient for tuning, we compare its selected
configuration with the choice from all five inner folds. Median agreement is 19.2\% for LightGBM and
17.8\% for the MLP across datasets, and is still incomplete when four folds are used. The one-fold
choices also score worse when evaluated across all five validation folds
(Appendix~Table~\ref{tab:inner-selection-stability}).
\FloatBarrier
